\section{서론}

\subsection{HNDL 위협과 PQC 전환 필요성}

양자 컴퓨터의 발전은 현재 인터넷 보안의 근간을 이루는 공개키 암호 체계에 실질적인 위협을 제기하고 있다.
Shor 알고리즘을 활용하면 RSA, ECC 등 현재 광범위하게 사용되는 공개키 알고리즘을 다항 시간 내에 해독할 수 있으며, 이는 양자 컴퓨터가 충분한 성능에 도달하는 시점에 기존 암호화 통신이 일괄적으로 무력화될 수 있음을 의미한다~[인용 추가].
더 시급한 문제는 ``지금 수집, 나중에 해독(Harvest Now, Decrypt Later, HNDL)'' 공격이다.
공격자는 현재 시점에서 TLS로 암호화된 네트워크 트래픽을 대량 수집해두고, 미래에 양자 컴퓨터를 이용해 복호화하는 방식으로 위협을 실현한다~[인용 추가].
TLS는 현재 인터넷에서 가장 광범위하게 사용되는 암호화 프로토콜이며, 그 키 교환 방식인 ECDHE는 Shor 알고리즘에 취약하다.
따라서 오늘 전송되는 TLS 트래픽이 이미 수집 대상이 될 수 있으며, PQC로의 전환은 미래의 과제가 아닌 현재의 과제다.

이러한 위협에 대응하여 NIST는 2024년 ML-KEM(FIPS~203), ML-DSA(FIPS~204), SLH-DSA(FIPS~205)를 PQC 표준으로 공식 확정하였다~[인용 추가].
표준화의 완료에 따라 각 조직은 기존 시스템의 TLS 구간을 PQC로 단계적으로 전환해야 하는 과제에 직면하였다.

\subsection{기존 TLS 전환 방식의 문제 정의}

단계적 PQC-TLS 전환은 알고리즘 교체만으로 완결되지 않는다.
실제 운영 환경에서는 다음 네 가지 문제가 구조적으로 발생한다.

첫째, TLS 설정 변경과 배포 과정이 분리되어 있어 설정 오류가 뒤늦게 발견된다.
설정 파일의 암호 정책이 올바른지 배포 이전에 자동으로 검증하는 체계가 없으면, 잘못된 설정이 운영 환경에 그대로 반영될 수 있다.
둘째, PQC 지원 클라이언트와 미지원 클라이언트가 공존하는 전환 과도기에 호환성 검증이 필요하다.
Hybrid PQC는 PQC 미지원 클라이언트에 대한 classical fallback을 제공하지만, 이 동작이 실제로 의도대로 이루어지는지 자동으로 확인하는 구조가 없다.
셋째, 전환 전후 사용 암호 자산을 추적하지 않으면 어떤 알고리즘이 제거·추가되었는지 증명하기 어렵다.
암호화 자산의 변화를 체계적으로 기록하고 비교하는 수단이 없으면 전환 진척도를 객관적으로 파악할 수 없다.
넷째, 성능 저하나 협상 실패 발생 시 자동 대응 구조가 없으면 운영 적용성이 떨어진다.
PQC 전환 과정에서는 기존 ECC 기반 구성과 다른 연산 및 통신 오버헤드가 발생할 수 있으며, 임계치를 초과하는 성능 저하가 발생할 경우 자동으로 이전 단계로 복구하는 메커니즘이 필요하다.

\subsection{연구 범위 및 위협 모델}

본 연구는 HNDL 공격에 직접 노출되는 TLS 통신 구간을 보호 대상으로 한다.
주요 위협은 현재 수집된 TLS 트래픽을 미래 양자 컴퓨터로 복호화하는 HNDL 공격이며, 전환 대상은 TLS 키 교환 알고리즘 및 인증서 서명 알고리즘을 중심으로 한다.

본 연구의 범위에서 제외되는 항목은 다음과 같다.
데이터베이스 암호화, 저장 데이터 암호화, 전체 애플리케이션 내부 암호 전환은 본 연구의 대상이 아니다.
Stage~3는 순수 PQC-only TLS의 표준화 및 구현체 지원이 아직 성숙하지 않은 점을 고려하여, Advanced Hybrid PQC (NIST L5) 구성으로 검증하며 순수 PQC-only TLS는 향후 연구로 남긴다.

\subsection{기여}

본 논문의 기여는 다음과 같다.
\begin{enumerate}[itemsep=4pt]
  \item \textbf{Policy as Code 기반 3단계 TLS 마이그레이션 모델 및 11-Step DevSecOps 파이프라인 설계·구현.}
  Classical ECC(Stage~1), Hybrid PQC(Stage~2), Advanced Hybrid PQC (NIST L5, Stage~3)의 단계적 전환 모델을 정의하고, 각 단계의 TLS 설정 변경·보안 스캔·동적 검증·성능 검증·자동 롤백을 통합한 11-Step 파이프라인을 구현하였다.

  \item \textbf{PQC 도구 체인 함정 6종 카탈로그화 및 자체 빌드·매트릭스 CI 기반 완화 구조 제시.}
  OQS 스택 통합 과정에서 발생하는 인프라 수준 함정 6종(F-1\textasciitilde F-6)을 트리거·증상·원인·방어책 단위로 구조화하였다.
  표준 SCA/SAST/SBOM 도구가 이를 사전 탐지하지 못함(0/6)을 정량 입증하고, 자체 빌드(5/6 예방)와 매트릭스 CI(6/6 감지)로 사각지대를 메우는 완화 구조를 제시한다.

  \item \textbf{CycloneDX 기반 CBOM Diff를 통한 암호 자산 전환 추적.}
  CycloneDX~1.7 기반 CBOM 생성 및 실행 이력 간 Diff를 구현하여, 전환 전후 암호 자산 변화를 BASELINE·IMPROVED·UNCHANGED·REGRESSED로 자동 판정한다.
  이를 통해 전환 진척도를 객관적으로 추적하고 회귀를 즉시 감지하는 구조를 제시한다.

  \item \textbf{LLM 기반 레거시 코드 마이그레이션 확장 모듈 구현 및 다층 검증 필요성 정량 평가.}
  Semgrep으로 레거시 암호를 탐지하고 LLM으로 ML-KEM/ML-DSA 기반 코드로 자동 변환한 뒤 PR을 자동 생성하는 확장 모듈을 구현하였다.
  80 trial 정량 평가를 통해 정적 AST 검증 단독으로는 런타임 실패를 차단할 수 없으며, 의미적 실행 검증이 필수적임을 확인하였다.
\end{enumerate}